\section{From source expressions to Lean proofs}
\label{sec:lean}

\subsection{A small semantic bridge rather than an unverified translation}

The first end-to-end target should be an origin-anchored inequality over a
literal exponential polynomial. Its proof construction has four distinct
objects:
\[
 \text{source expression}
 \xrightarrow{\text{proved equality}}
 \operatorname{eval}(p)
 \xrightarrow{\text{checked ladder}}
 \operatorname{NonnegativeRay}(p,a)
 \xrightarrow{\text{domain transport}}
 \text{original goal}.
\]
The arrows represent proofs, not merely successful parsing or a matching hash.
The final arrow is trivial only when the source and certificate domains and
strictness already coincide.

A proposed Lean interface is:
\par\noindent\begin{minipage}{\linewidth}
\begin{lstlisting}
-- Interface sketch; names are proposed, not existing Forge declarations.
structure ExpPoly where
  blocks : Array (Rat x Array Rat)

structure ReifiedExpPoly where
  source      : Expr
  polynomial  : ExpPoly
  evalEq      : Expr  -- proof: source = eval polynomial at the chosen variable
  variable    : FVarId
  contextKey  : ContextKey

structure LadderReceipt where
  anchor  : Rat
  factors : Array Rat
  pointCertificates : Array PointCertificate

checkLadder_sound :
  checkLadder p c = true ->
  forall x : Real, c.anchor <= x -> 0 <= eval p x
\end{lstlisting}
\end{minipage}\par

The syntax in this sketch uses readable placeholders such as \code{x} for a
product and \code{Real} for the real type; it is not advertised as compiling
Lean. The supplied \code{LadderBridge.lean}, in contrast, is actual candidate
Lean source with the full intended theorem and proof, but remains uncompiled.

The internal polynomial representation should have a well-formedness predicate
for ordered unique rates, trimmed coefficient arrays, and exact rationals.
Either use a subtype carrying that proof, or canonicalise before every checked
operation. Do not permit two distinct encodings of the same block to bypass a
terminal-zero or source-equality check.

\subsection{Logarithmic charts expand the useful source language}

For $x>0$, put $t=\log x$, so $x=e^t$. An expression
\[
 \sum_{j,k}c_{j,k}x^{r_j}(\log x)^k
\]
with rational exponents becomes
$\sum_{j,k}c_{j,k}t^ke^{r_jt}$ once the real-power identities and their domain
conditions have been proved. Integer negative powers also fit this pattern.
The map is particularly convenient on $x\ge1$, which corresponds to $t\ge0$.
On $0<x\le1$, use $u=-\log x\ge0$ and reflection.

The chart is not permission to rewrite real powers at nonpositive arguments.
A denominator must be proved positive or nonzero before clearing it; a square
root must be connected to its exponential representation on the positive
domain; and $\log1=0$ must preserve whether the original endpoint is included.

\begin{example}[Entropy inequality]
For $x>0$,
\[
 x\log x-x+1\ge0.
\]
Its chart is $f(t)=(t-1)e^t+1$. On $t\ge0$, native factors $(0,1,1)$ have
initial vector $(0,0,1)$. On $t\le0$, reflect: $f(-u)=1-(u+1)e^{-u}$,
whose factors $(-1,-1,0)$ also have initial vector $(0,0,1)$.
Thus the source inequality follows on both sides of one, with equality at one
and strict positivity elsewhere.
\end{example}

\begin{example}[A rational lower bound for the logarithm]
For $x\ge1$,
\[
 \log x\ge\frac{2(x-1)}{x+1}.
\]
After proving $x+1>0$ and writing $x=e^t$, the target is
\[
 (t-2)e^t+t+2\ge0\qquad(t\ge0).
\]
Factors $(0,0,1,1)$ have initial vector $(0,0,0,1)$. The third-order boundary
contact is handled exactly; no interval precision is spent resolving it.
\end{example}

\begin{example}[A square-root upper bound for the logarithm]
For $x\ge1$,
\[
 \log x\le\frac{x-1}{\sqrt x}.
\]
The proved positive-domain chart gives
\[
 e^{t/2}-e^{-t/2}-t\ge0\qquad(t\ge0).
\]
Factors $(-1/2,0,0,1/2)$ have initial vector $(0,0,0,1/4)$.
\end{example}

The archive records certificates for these \emph{chart polynomials}. It does
not contain an executed Lean reifier proving the original logarithmic source
goals. The transformations above are mathematical derivations and concrete
acceptance targets for that reifier; they are not additional kernel-checked
experiment rows.

\subsection{Composing the worker with existing algebra}

Several useful extensions require no new analytic checker. For a quotient
$f/g$ with a proved nonzero denominator, its derivative has numerator
$f'g-fg'$ and denominator $g^2>0$. The numerator is again an exponential
polynomial if $f$ and $g$ are. Certifying its sign establishes monotonicity;
certifying the corresponding next derivative numerator can establish convexity.
The denominator obligations must remain attached to the original domain.

A common positive factor may be removed. For example,
\[
 e^x\ge e^y(1+x-y)
\]
reduces, after multiplication by $e^{-y}>0$, to
$e^t\ge1+t$ with $t=x-y$. A proof over all real $t$ combines the two reflected
ray certificates. This handles a multivariate source through a one-dimensional
proved reduction; it is not a general multivariate exponential procedure.

Similarly, polynomial factors known to be nonnegative should be retained
outside the ladder whenever possible. Expanding a square only to force an
analytic search to rediscover positivity wastes both algebraic structure and
proof size. The router should compare the expanded derivative representation
with a factored source view rather than making either view globally canonical.

\subsection{Concrete libraries to reuse}

\paragraph{Mathlib derivative and monotonicity theorems.}
Use the explicit-derivative API in
\pathline{Mathlib/Analysis/Calculus/Deriv/MeanValue.lean}
and
\pathline{Mathlib/Analysis/SpecialFunctions/ExpDeriv.lean}.
These expose the monotonicity and exponential chain-rule ingredients for the
step theorem~\cite{mathlib-meanvalue,mathlib-expderiv}. Explicit
\code{HasDerivAt} witnesses avoid the default-zero behaviour of a totalised
\code{deriv} when differentiability has not been established.

\paragraph{LeanCert for bounded analytic obligations.}
LeanCert's maintained interface documents checked interval evaluation,
automatic differentiation, bounds, optimisation, and root certificates.
Programmatic entry points include \code{LeanCert.evalInterval},
\code{LeanCert.evalInterval_correct}, \code{LeanCert.API.Bounds}, and
\code{LeanCert.API.Optimization}. Its documented automatic tail fragment is
currently reciprocal-power bounds over naturals, rather than general
exponential or logarithmic tails~\cite{leancert}. Thus it is a concrete
candidate for the compact-domain and point-bound adapter, not a reason to
reimplement its interval evaluator inside Forge.

Trust selection must be explicit. LeanCert documents a native mode using
\code{native_decide}, a kernel mode, and an automatic mode that can fall back.
For the strict kernel-only route envisioned here, request
\code{leancert (trust := kernel)} or use the corresponding checked API without
native fallback. This archive inspected those interfaces; it did not install,
run, or benchmark LeanCert.

\paragraph{IntervalBounds for small concrete endpoint estimates.}
The \code{peti12352/lean-interval-bounds} project provides \code{approx},
\code{interval_decide}, and \code{norm_num} extensions for exact transcendental
bounds. Its documented approach avoids floating-point proof arithmetic and
\code{native_decide}. It is narrower than a quantified interval worker, but
well matched to rational-anchor estimates. Its README reports a
\code{v4.33.0-rc1} toolchain, different from the audited Forge baseline;
compatibility must be compiled rather than assumed~\cite{intervalbounds}.

\paragraph{Exact arithmetic and algebraic side conditions.}
The prototype needs only Python's \code{Fraction}, integer arithmetic, arrays,
and maps. A first Lean implementation can use rational coefficients and proved
sparse operations, plus ordinary ring normalisation and numerical proofs for
small emitted identities. No SAT, SMT, SDP, or numerical eigenvalue library is
required for the native ladder or the total tail constructor. Adding one would
complicate deployment without solving a missing subproblem in these algorithms.

The library recommendation is therefore narrow: Mathlib for the analytic bridge,
a proved exact polynomial representation for the certificate arithmetic,
LeanCert or IntervalBounds for endpoint and compact obligations, and Forge's
existing algebraic worker for factors. These pieces have separate responsibilities
and no implied claim of plug-and-play version compatibility.

\subsection{Candidate proof emission in the archive}

The supplied \code{LadderBridge.lean} states the one-step theorem in terms of
arbitrary real functions and explicit derivative proofs. The emitter then
creates a private definition for every $q_j$, proves its derivative identity,
proves its initial value by exact arithmetic, and applies the step theorem
backwards from the zero function. This is intentionally a proof-producing
path, not an invocation of the Python verifier from Lean.

The four generated modules cover the Taylor remainder of order four, the
positive entropy chart, the rational logarithm-bound chart, and the
auxiliary-factor example with $\varepsilon=1$. Their final statements
say nonnegativity on the closed nonnegative ray. Some corresponding Python
receipts establish stronger strict conclusions, but the emitted Lean statements
are not presented as closing those stronger requests.

Every candidate requests an axiom inventory. None of those requests was
executed because this environment has no Lean executable. The first integration
milestone must compile and repair the bridge and all four replays without
weakening their statements or introducing proof holes. Only then should proof
emission be replaced with a reflected checker for larger certificates.

\subsection{Why a reflected checker is the second step, not the first}

A reflected checker should prove soundness once for evaluation, shifted
 differentiation, point certificates, and ladder recursion. It then reduces a
closed Boolean certificate check and invokes the soundness theorem. This can
avoid a large tower of repeated ring proofs, especially for auxiliary-factor
certificates with long integer coefficients.

However, reflection introduces additional obligations: correctness of rational
casts, polynomial evaluation, sorting and trimming, and the executable
arithmetic backend. A small emitted proof is an excellent reference oracle
while those lemmas are developed. Keep the source-emission path as a fallback
and differential test even after reflection becomes the production route.
